\begin{graphillustration}{Пример: подграф может опустить ребро $ac$, а порождённый подграф $G[W]$ обязан его содержать.}{fig:notation-subgraphs}
\begin{tikzpicture}[font=\small]
\begin{scope}[shift={(0,0)}]
\node[gi vertex] (a) at (0,0) {$a$};
\node[gi vertex] (b) at (1.3,1) {$b$};
\node[gi vertex] (c) at (2.6,0) {$c$};
\node[gi vertex] (d) at (1.3,-1) {$d$};
\draw[gi edge] (a) -- (b);
\draw[gi edge] (b) -- (c);
\draw[gi edge] (c) -- (a);
\draw[gi edge] (c) -- (d);
\draw[gi edge] (d) -- (a);
\node[gi note] at (1.3,-1.6) {$G$};
\end{scope}
\begin{scope}[shift={(4.5,0)}]
\node[gi vertex] (a) at (0,0) {$a$};
\node[gi vertex] (b) at (1.3,1) {$b$};
\node[gi vertex] (c) at (2.6,0) {$c$};
\draw[gi edge] (a) -- (b);
\draw[gi edge] (b) -- (c);
\node[gi note] at (1.3,-1.6) {$H\subset G$, $W=\{a,b,c\}$};
\end{scope}
\begin{scope}[shift={(9,0)}]
\node[gi vertex] (a) at (0,0) {$a$};
\node[gi vertex] (b) at (1.3,1) {$b$};
\node[gi vertex] (c) at (2.6,0) {$c$};
\draw[gi edge] (a) -- (b);
\draw[gi edge] (b) -- (c);
\draw[gi edge] (c) -- (a);
\node[gi note] at (1.3,-1.6) {$G[W]$};
\end{scope}
\end{tikzpicture}
\end{graphillustration}
